\documentclass[10pt,a4paper]{article}
\usepackage[spanish,es-nodecimaldot]{babel}
\usepackage[utf8]{inputenc}
\usepackage[T1]{fontenc}
\usepackage{lmodern}
\usepackage[top=1.1cm,bottom=1.1cm,left=1.5cm,right=1.5cm]{geometry}
\usepackage{amsmath,amssymb}
\usepackage[table]{xcolor}
\usepackage{enumitem}
\usepackage[normalem]{ulem}
\usepackage{tikz}

\setlength{\parindent}{0pt}
\setlength{\parskip}{2pt}
\pagestyle{empty}

\AtBeginDocument{%
  \setlength{\abovedisplayskip}{3pt}%
  \setlength{\belowdisplayskip}{3pt}%
  \setlength{\abovedisplayshortskip}{2pt}%
  \setlength{\belowdisplayshortskip}{2pt}%
}

\begin{document}
\shorthandoff{<>}

\begin{center}
{\large\bfseries Metodos de los minimos cuadrados}
\end{center}
\vspace{-2pt}

Se requiere analizar la \textbf{relación} entre dos \textbf{variables cuantitativas}, para determinar si están asociadas, y si una puede predecir el valor de otra.
\textbf{El coefic. de corr. no sirve:} no explica una variable respecto de otra

\textbf{\uline{Regresion lineal}}

1) Tenemos un diagrama de dispersión donde hay dos variables relacionadas:
\begin{itemize}[leftmargin=16pt,itemsep=0pt,topsep=1pt,parsep=0pt]
  \item La independiente o \textbf{explicativa ``x''} \quad ej: edad
  \item La dependiente o \textbf{explicada ``y''} \qquad ej: tensión arterial
\end{itemize}

2) Queremos encontrar la recta que mejor ajuste los valores, y pueda ser utilizada para predecir los valores:
\begin{itemize}[leftmargin=16pt,itemsep=0pt,topsep=1pt,parsep=0pt]
  \item Para esto usamos el \textbf{método de los mínimos cuadrados}, que elige la recta de regresión que minimiza los errores de las diferencias vertices.
\end{itemize}

\begin{minipage}[t]{0.55\textwidth}
\vspace{0pt}
3) La recta que mejor se ajusta es: $y = \alpha + \beta.x$, que será estimada por:
\[
\hat{y} = a + b\cdot x
\]
{\small a y b (parámetros muestrales) son estimadores de $\alpha$ y $\beta$ (parámetros poblacionales)}
\[
e_i = y_i - \hat{y}_i
\qquad
\text{\small diferencias verticales} \leftarrow \text{\small (regresión y sobre x)}
\]
\end{minipage}\hfill
\begin{minipage}[t]{0.42\textwidth}
\vspace{2pt}
\begin{center}
\begin{tikzpicture}[scale=0.62, every node/.style={font=\scriptsize}]
  \draw[->] (0,0) -- (7.2,0) node[right] {$x$};
  \draw[->] (0,0) -- (0,5.2) node[above] {$y$};
  \draw[thick, black] (0.5,1.0) -- (6.8,4.6) node[right] {$y=\alpha+\beta.x$};
  \foreach \p in {(1.1,1.3),(1.6,1.9),(2.3,2.05),(3.0,2.55),(3.55,2.7),(4.2,3.15)}
    \fill[red] \p circle (1.3pt);
  \draw[dashed, blue] (3.55,0) -- (3.55,2.85) node[below, black, xshift=-2pt, yshift=-11pt] {$x_i$};
  \draw[dashed] (3.55,2.85) -- (0,2.85) node[left] {$\hat{y}$};
  \draw[<->, thick, violet] (3.55,2.7) -- (3.55,3.05);
  \node[violet, left] at (3.4,2.9) {$\varepsilon_i$};
  \node[above] at (5.6,1.6) {\small $y$};
\end{tikzpicture}
\end{center}
\vspace{-2pt}
\centerline{\footnotesize No es viable la suma de los errores,}
\centerline{\footnotesize pq se cancelan los (+) con los (-)}
\end{minipage}

4) Se busca \textbf{minimizar la suma de los cuadrados de las diferencias} entre los valores reales y los valores estimados $\Rightarrow$ a y b se eligirán con:
\[
\sum_{i=1}^{n} \left[ y_i - (a+b\cdot x_i) \right]^2
\qquad
\begin{minipage}{0.5\textwidth}
{\footnotesize Al cuadrado para que los errores no se cancelen, no utilizamos el valor absoluto debido a un problema de amiguedad en el $f(x)' = 0$}
\end{minipage}
\]

5) Para minimizar la suma debemos hacer las \textbf{derivadas parciales} de los coeficientes a y b, e igualarlos a 0. Esto es debido a que la función es convexa entonces al igualar a 0 se encuentra el mínimo de la función
\[
\sum_{i=1}^{n} y_i = a\cdot n + b\sum_{i=1}^{n} x_i
\qquad\qquad
\sum_{i=1}^{n} y_i\cdot x_i = a\sum_{i=1}^{n} x_i + b\sum_{i=1}^{n} (x_i)^2
\]

6) Formamos un Sist de Ecuaciones Normales con las derivadas:

\begin{minipage}[t]{0.52\textwidth}
\vspace{0pt}
\[
\begin{bmatrix}
n & \displaystyle\sum_{i=1}^{n} x_i \\[6pt]
\displaystyle\sum_{i=1}^{n} x_i & \displaystyle\sum_{i=1}^{n} (x_i)^2
\end{bmatrix}
\begin{pmatrix} a \\ b \end{pmatrix}
=
\begin{bmatrix}
\displaystyle\sum_{i=1}^{n} y_i \\[6pt]
\displaystyle\sum_{i=1}^{n} (y_i\cdot x_i)
\end{bmatrix}
\]
\end{minipage}\hfill
\begin{minipage}[t]{0.44\textwidth}
\vspace{4pt}
{\small
Resolviendo el SEN por determinantes \textbf{obtenemos los valores de a y b} que mejor se ajustan a la recta

Se puede calcular a con \texttt{intercept(x,y)} y b con \texttt{slope(x,y)}
}
\end{minipage}

\end{document}
